\section{RISC-V CoVE-IO / TEE-I/O}

\begin{frame}[t,shrink=6]{RISC-V CoVE-IO / TEE-I/O：方向开场}
\DeckKicker{方向开场}
\SlideMeta{证据等级}{方向综述}{来源状态：公开材料综合}
{\large\bfseries\color{Ink} RISC-V CoVE-IO / TEE-I/O 的核心是先界定保护对象、管理边界和证据强度，再用三篇主讲材料串起技术演进。\par}\vspace{1.8mm}
\begin{columns}[T,totalwidth=\textwidth]
  \begin{column}{0.45\textwidth}
    \fcolorbox{Rule}{Paper}{%
  \begin{minipage}[t]{0.97\linewidth}
    \vspace{1.1mm}
    \hspace{1.4mm}\begin{minipage}{0.92\linewidth}
      \PanelTitle{内容说明}
      \scriptsize \begin{itemize}
  \item 把设备身份、DMA/MMIO、interrupt、link security 与 I/O translation 放到同一边界。
  \item 固定三篇主讲材料；`riscv\_aia\_2023` 作为 trusted interrupt auxiliary，不计入三篇主讲。
  \item 先讲方向痛点和设计轴，再逐篇说明贡献、限制、证据强度和可落地条件。
  \item 对规范、Survey、draft、vendor evidence 保持显式标注，避免把背景材料写成一手系统证明。
\end{itemize}
    \end{minipage}
    \vspace{1mm}
  \end{minipage}}
  \end{column}
  \begin{column}{0.49\textwidth}
    \fcolorbox{Rule}{Panel}{%
  \begin{minipage}[t]{0.97\linewidth}
    \vspace{1.1mm}
    \hspace{1.4mm}\begin{minipage}{0.92\linewidth}
      \PanelTitle{三篇主讲材料的分工}
      \scriptsize {\tiny
\begin{tabularx}{\linewidth}{@{}YYY@{}}
\toprule
\textbf{论文} & \textbf{定位} & \textbf{证据边界} \\
\midrule
sIOPMP: scalable I/O protection for TEE & 开创/基础入口 & E1 / 同行评审改进证据 \\
RISC-V CoVE-IO 规范草案 & 代表性改进 & E3 / 草案/未批准规范 \\
RISC-V IOMMU 支撑 CoVE-IO & 当前 SOTA/补充边界 & E0 / Spec-standard SOTA \\
\bottomrule
\end{tabularx}
}
    \end{minipage}
    \vspace{1mm}
  \end{minipage}}
  \end{column}
\end{columns}
\vspace{1mm}
{\tiny\textcolor{Muted}{引用依据：本页综合三篇主讲材料的研究定位、证据等级和公开来源状态。}}
\end{frame}


\begin{frame}[t,shrink=6]{sIOPMP: scalable I/O protection for TEE：内容摘要}
\DeckKicker{内容摘要}
\SlideMeta{证据等级}{E1 / 同行评审改进证据}{来源状态：本地 PDF 已核验}
{\large\bfseries\color{Ink} 提出面向 TEE 的 scalable I/O protection，聚焦设备 access control 和 metadata/checkin。\par}\vspace{1.8mm}
\begin{columns}[T,totalwidth=\textwidth]
  \begin{column}{0.45\textwidth}
    \fcolorbox{Rule}{Paper}{%
  \begin{minipage}[t]{0.97\linewidth}
    \vspace{1.1mm}
    \hspace{1.4mm}\begin{minipage}{0.92\linewidth}
      \PanelTitle{内容说明}
      \scriptsize \begin{itemize}
  \item 提出面向 TEE 的 scalable I/O protection，聚焦设备 access control 和 metadata/checking 可扩展性
  \item sIOPMP is a peer-reviewed mechanism starting point for RISC-V trusted I/O and DMA protectio。
\end{itemize}
    \end{minipage}
    \vspace{1mm}
  \end{minipage}}
  \end{column}
  \begin{column}{0.49\textwidth}
    \fcolorbox{Rule}{Panel}{%
  \begin{minipage}[t]{0.97\linewidth}
    \vspace{1.1mm}
    \hspace{1.4mm}\begin{minipage}{0.92\linewidth}
      \PanelTitle{证据定位}
      \scriptsize {\tiny
\begin{tabularx}{\linewidth}{@{}YY@{}}
\toprule
\textbf{字段} & \textbf{内容} \\
\midrule
论文定位 & 开创/基础入口 \\
材料类型 & 系统论文 \\
证据等级 & E1 / 同行评审改进证据 \\
来源状态 & 本地 PDF 已核验 \\
\bottomrule
\end{tabularx}
}
    \end{minipage}
    \vspace{1mm}
  \end{minipage}}
  \end{column}
\end{columns}
\vspace{1mm}
{\tiny\textcolor{Muted}{引用依据：引用依据为论文原文、官方规范或公开材料；本页结论按 E1 / 同行评审改进证据 的证据等级使用，不外推到未验证平台。}}
\end{frame}


\begin{frame}[t,shrink=6]{sIOPMP: scalable I/O protection for TEE：研究背景}
\DeckKicker{研究背景}
\SlideMeta{证据等级}{E1 / 同行评审改进证据}{来源状态：本地 PDF 已核验}
{\large\bfseries\color{Ink} DMA 设备可绕过 CPU memory isolation，I/O 保护是 TEE 的关键缺口\par}\vspace{1.8mm}
\begin{columns}[T,totalwidth=\textwidth]
  \begin{column}{0.45\textwidth}
    \fcolorbox{Rule}{Paper}{%
  \begin{minipage}[t]{0.97\linewidth}
    \vspace{1.1mm}
    \hspace{1.4mm}\begin{minipage}{0.92\linewidth}
      \PanelTitle{内容说明}
      \scriptsize \begin{itemize}
  \item DMA 设备可绕过 CPU memory isolation，I/O 保护是 TEE 的关键缺口
  \item 只保护 CPU 页表或 PMP 不足以覆盖设备侧访问
\end{itemize}
    \end{minipage}
    \vspace{1mm}
  \end{minipage}}
  \end{column}
  \begin{column}{0.49\textwidth}
    \fcolorbox{Rule}{Panel}{%
  \begin{minipage}[t]{0.97\linewidth}
    \vspace{1.1mm}
    \hspace{1.4mm}\begin{minipage}{0.92\linewidth}
      \PanelTitle{问题边界}
      \scriptsize {\tiny
\begin{tabularx}{\linewidth}{@{}YY@{}}
\toprule
\textbf{维度} & \textbf{说明} \\
\midrule
保护对象 & RISC-V CoVE-IO / TEE-I/O \\
传统瓶颈 & 把设备身份、DMA/MMIO、interrupt、link security 与 I/O translation 放到同一边界。 \\
本文入口 & sIOPMP: scalable I/O protection for TEE \\
证据限制 & E1 / 同行评审改进证据 \\
\bottomrule
\end{tabularx}
}
    \end{minipage}
    \vspace{1mm}
  \end{minipage}}
  \end{column}
\end{columns}
\vspace{1mm}
{\tiny\textcolor{Muted}{引用依据：引用依据为论文原文、官方规范或公开材料；本页结论按 E1 / 同行评审改进证据 的证据等级使用，不外推到未验证平台。}}
\end{frame}


\begin{frame}[t,shrink=6]{sIOPMP: scalable I/O protection for TEE：解决方案}
\DeckKicker{解决方案}
\SlideMeta{证据等级}{E1 / 同行评审改进证据}{来源状态：本地 PDF 已核验}
{\large\bfseries\color{Ink} 设计 scalable I/O protection metadata/checking 机制，提高设备访问控制效率并支撑更大规模 I/O pro。\par}\vspace{1.8mm}
\begin{columns}[T,totalwidth=\textwidth]
  \begin{column}{0.45\textwidth}
    \fcolorbox{Rule}{Paper}{%
  \begin{minipage}[t]{0.97\linewidth}
    \vspace{1.1mm}
    \hspace{1.4mm}\begin{minipage}{0.92\linewidth}
      \PanelTitle{内容说明}
      \scriptsize \begin{itemize}
  \item 设计 scalable I/O protection metadata/checking 机制，提高设备访问控制效率并支撑更大规模 I/O protection
\end{itemize}
    \end{minipage}
    \vspace{1mm}
  \end{minipage}}
  \end{column}
  \begin{column}{0.49\textwidth}
    \fcolorbox{Rule}{Panel}{%
  \begin{minipage}[t]{0.97\linewidth}
    \vspace{1.1mm}
    \hspace{1.4mm}\begin{minipage}{0.92\linewidth}
      \PanelTitle{核心设计拆解}
      \scriptsize \begin{itemize}
  \item 01. Scalable I/O protection metadata
  \item 02. Device access-control checking path
  \item 03. TEE-oriented DMA isolation substrate
\end{itemize}
    \end{minipage}
    \vspace{1mm}
  \end{minipage}}
  \end{column}
\end{columns}
\vspace{1mm}
{\tiny\textcolor{Muted}{引用依据：引用依据为论文原文、官方规范或公开材料；本页结论按 E1 / 同行评审改进证据 的证据等级使用，不外推到未验证平台。}}
\end{frame}


\begin{frame}[t,shrink=6]{sIOPMP: scalable I/O protection for TEE：实验结果}
\DeckKicker{实验结果}
\SlideMeta{证据等级}{E1 / 同行评审改进证据}{来源状态：本地 PDF 已核验}
{\large\bfseries\color{Ink} ASPLOS 论文给出吞吐和开销评估\par}\vspace{1.8mm}
\begin{columns}[T,totalwidth=\textwidth]
  \begin{column}{0.45\textwidth}
    \fcolorbox{Rule}{Paper}{%
  \begin{minipage}[t]{0.97\linewidth}
    \vspace{1.1mm}
    \hspace{1.4mm}\begin{minipage}{0.92\linewidth}
      \PanelTitle{内容说明}
      \scriptsize \begin{itemize}
  \item ASPLOS 论文给出吞吐和开销评估
  \item 具体数值以本地 PDF 为准
\end{itemize}
    \end{minipage}
    \vspace{1mm}
  \end{minipage}}
  \end{column}
  \begin{column}{0.49\textwidth}
    \fcolorbox{Rule}{Panel}{%
  \begin{minipage}[t]{0.97\linewidth}
    \vspace{1.1mm}
    \hspace{1.4mm}\begin{minipage}{0.92\linewidth}
      \PanelTitle{实验/证据边界}
      \scriptsize {\tiny
\begin{tabularx}{\linewidth}{@{}YYY@{}}
\toprule
\textbf{对象} & \textbf{可支撑} & \textbf{边界} \\
\midrule
数据/来源 & sIOPMP: scalable I/O protection for TEE & 本地 PDF 已核验 \\
Baseline/对照 & 以论文或规范原文为准 & 不跨材料外推 \\
指标/对象 & 只使用当前来源可证明的信息 & E1 / 同行评审改进证据 \\
使用边界 & 可进入本方向演进图 & 不能替代更强证据 \\
\bottomrule
\end{tabularx}
}
    \end{minipage}
    \vspace{1mm}
  \end{minipage}}
  \end{column}
\end{columns}
\vspace{1mm}
{\tiny\textcolor{Muted}{引用依据：引用依据为论文原文、官方规范或公开材料；本页结论按 E1 / 同行评审改进证据 的证据等级使用，不外推到未验证平台。}}
\end{frame}


\begin{frame}[t,shrink=6]{sIOPMP: scalable I/O protection for TEE：文章评价}
\DeckKicker{文章评价}
\SlideMeta{证据等级}{E1 / 同行评审改进证据}{来源状态：本地 PDF 已核验}
{\large\bfseries\color{Ink} RISC-V I/O protection 的 peer-reviewed SOTA，可作为 CoVE-IO access-control 前置材料\par}\vspace{1.8mm}
\begin{columns}[T,totalwidth=\textwidth]
  \begin{column}{0.45\textwidth}
    \fcolorbox{Rule}{Paper}{%
  \begin{minipage}[t]{0.97\linewidth}
    \vspace{1.1mm}
    \hspace{1.4mm}\begin{minipage}{0.92\linewidth}
      \PanelTitle{内容说明}
      \scriptsize \begin{itemize}
  \item 设计优势：RISC-V I/O protection 的 peer-reviewed SOTA，可作为 CoVE-IO access-control 前置材料。
  \item 局限性：主要解决 access control，不覆盖完整 device identity、TDISP、link security 或 trusted interrupt pat。
  \item 商业化潜力：对高性能设备直通有价值；风险在硬件改造、IOMMU/IOPMP 组合和平台可观测性。
  \item 证据边界：E1 / Peer-reviewed SOTA；本地 PDF 已核验。
\end{itemize}
    \end{minipage}
    \vspace{1mm}
  \end{minipage}}
  \end{column}
  \begin{column}{0.49\textwidth}
    \fcolorbox{Rule}{Panel}{%
  \begin{minipage}[t]{0.97\linewidth}
    \vspace{1.1mm}
    \hspace{1.4mm}\begin{minipage}{0.92\linewidth}
      \PanelTitle{文章评价}
      \scriptsize {\tiny
\begin{tabularx}{\linewidth}{@{}YY@{}}
\toprule
\textbf{维度} & \textbf{判断} \\
\midrule
设计优势 & RISC-V I/O protection 的 peer-reviewed SOTA，可作为 CoVE-IO access-control 前置材料。 \\
主要局限 & 主要解决 access control，不覆盖完整 device identity、TDISP、link security 或 trusted interru。 \\
商业落地 & 对高性能设备直通有价值；风险在硬件改造、IOMMU/IOPMP 组合和平台可观测性。 \\
讲稿定位 & 开创/基础入口; 证据强度 E1 \\
\bottomrule
\end{tabularx}
}
    \end{minipage}
    \vspace{1mm}
  \end{minipage}}
  \end{column}
\end{columns}
\vspace{1mm}
{\tiny\textcolor{Muted}{引用依据：引用依据为论文原文、官方规范或公开材料；本页结论按 E1 / 同行评审改进证据 的证据等级使用，不外推到未验证平台。}}
\end{frame}


\begin{frame}[t,shrink=6]{RISC-V CoVE-IO 规范草案：内容摘要}
\DeckKicker{内容摘要}
\SlideMeta{证据等级}{E3 / 草案/未批准规范}{来源状态：本地 PDF 已核验}
{\large\bfseries\color{Ink} 定义 RISC-V CoVE-IO confidential I/O 架构，把 TVM 设备使用扩展到 device identity、DMA/M。\par}\vspace{1.8mm}
\begin{columns}[T,totalwidth=\textwidth]
  \begin{column}{0.45\textwidth}
    \fcolorbox{Rule}{Paper}{%
  \begin{minipage}[t]{0.97\linewidth}
    \vspace{1.1mm}
    \hspace{1.4mm}\begin{minipage}{0.92\linewidth}
      \PanelTitle{内容说明}
      \scriptsize \begin{itemize}
  \item 定义 RISC-V CoVE-IO confidential I/O 架构，把 TVM 设备使用扩展到 device identity、DMA/MMIO、interrupt 和 li。
  \item CoVE-IO and TEE-I/O draft is the current standards-track SOTA for RISC-V confidential I/O.
\end{itemize}
    \end{minipage}
    \vspace{1mm}
  \end{minipage}}
  \end{column}
  \begin{column}{0.49\textwidth}
    \fcolorbox{Rule}{Panel}{%
  \begin{minipage}[t]{0.97\linewidth}
    \vspace{1.1mm}
    \hspace{1.4mm}\begin{minipage}{0.92\linewidth}
      \PanelTitle{证据定位}
      \scriptsize {\tiny
\begin{tabularx}{\linewidth}{@{}YY@{}}
\toprule
\textbf{字段} & \textbf{内容} \\
\midrule
论文定位 & 代表性改进 \\
材料类型 & 规范/标准材料 \\
证据等级 & E3 / 草案/未批准规范 \\
来源状态 & 本地 PDF 已核验 \\
\bottomrule
\end{tabularx}
}
    \end{minipage}
    \vspace{1mm}
  \end{minipage}}
  \end{column}
\end{columns}
\vspace{1mm}
{\tiny\textcolor{Muted}{引用依据：引用依据为论文原文、官方规范或公开材料；本页结论按 E3 / 草案/未批准规范 的证据等级使用，不外推到未验证平台。}}
\end{frame}


\begin{frame}[t,shrink=6]{RISC-V CoVE-IO 规范草案：研究背景}
\DeckKicker{研究背景}
\SlideMeta{证据等级}{E3 / 草案/未批准规范}{来源状态：本地 PDF 已核验}
{\large\bfseries\color{Ink} TVM 使用真实设备时，安全边界必须覆盖设备身份、DMA/MMIO、interrupt delivery、link protection 和管理。\par}\vspace{1.8mm}
\begin{columns}[T,totalwidth=\textwidth]
  \begin{column}{0.45\textwidth}
    \fcolorbox{Rule}{Paper}{%
  \begin{minipage}[t]{0.97\linewidth}
    \vspace{1.1mm}
    \hspace{1.4mm}\begin{minipage}{0.92\linewidth}
      \PanelTitle{内容说明}
      \scriptsize \begin{itemize}
  \item TVM 使用真实设备时，安全边界必须覆盖设备身份、DMA/MMIO、interrupt delivery、link protection 和管理 control plane
\end{itemize}
    \end{minipage}
    \vspace{1mm}
  \end{minipage}}
  \end{column}
  \begin{column}{0.49\textwidth}
    \fcolorbox{Rule}{Panel}{%
  \begin{minipage}[t]{0.97\linewidth}
    \vspace{1.1mm}
    \hspace{1.4mm}\begin{minipage}{0.92\linewidth}
      \PanelTitle{问题边界}
      \scriptsize {\tiny
\begin{tabularx}{\linewidth}{@{}YY@{}}
\toprule
\textbf{维度} & \textbf{说明} \\
\midrule
保护对象 & RISC-V CoVE-IO / TEE-I/O \\
传统瓶颈 & 把设备身份、DMA/MMIO、interrupt、link security 与 I/O translation 放到同一边界。 \\
本文入口 & RISC-V CoVE-IO 规范草案 \\
证据限制 & E3 / 草案/未批准规范 \\
\bottomrule
\end{tabularx}
}
    \end{minipage}
    \vspace{1mm}
  \end{minipage}}
  \end{column}
\end{columns}
\vspace{1mm}
{\tiny\textcolor{Muted}{引用依据：引用依据为论文原文、官方规范或公开材料；本页结论按 E3 / 草案/未批准规范 的证据等级使用，不外推到未验证平台。}}
\end{frame}


\begin{frame}[t,shrink=6]{RISC-V CoVE-IO 规范草案：解决方案}
\DeckKicker{解决方案}
\SlideMeta{证据等级}{E3 / 草案/未批准规范}{来源状态：本地 PDF 已核验}
{\large\bfseries\color{Ink} 用 TDI/TDM/DSM、SPDM、TDISP、PCIe IDE、trusted MSI 等机制组合描述 draft confidential。\par}\vspace{1.8mm}
\begin{columns}[T,totalwidth=\textwidth]
  \begin{column}{0.45\textwidth}
    \fcolorbox{Rule}{Paper}{%
  \begin{minipage}[t]{0.97\linewidth}
    \vspace{1.1mm}
    \hspace{1.4mm}\begin{minipage}{0.92\linewidth}
      \PanelTitle{内容说明}
      \scriptsize \begin{itemize}
  \item 用 TDI/TDM/DSM、SPDM、TDISP、PCIe IDE、trusted MSI 等机制组合描述 draft confidential I/O lifecycle
\end{itemize}
    \end{minipage}
    \vspace{1mm}
  \end{minipage}}
  \end{column}
  \begin{column}{0.49\textwidth}
    \fcolorbox{Rule}{Panel}{%
  \begin{minipage}[t]{0.97\linewidth}
    \vspace{1.1mm}
    \hspace{1.4mm}\begin{minipage}{0.92\linewidth}
      \PanelTitle{核心设计拆解}
      \scriptsize \begin{itemize}
  \item 01. TDI/TDM/DSM device-domain vocabulary
  \item 02. SPDM/TDISP and PCIe IDE linkage
  \item 03. Trusted MSI and interrupt-path concepts
\end{itemize}
    \end{minipage}
    \vspace{1mm}
  \end{minipage}}
  \end{column}
\end{columns}
\vspace{1mm}
{\tiny\textcolor{Muted}{引用依据：引用依据为论文原文、官方规范或公开材料；本页结论按 E3 / 草案/未批准规范 的证据等级使用，不外推到未验证平台。}}
\end{frame}


\begin{frame}[t,shrink=6]{RISC-V CoVE-IO 规范草案：实验结果}
\DeckKicker{实验结果}
\SlideMeta{证据等级}{E3 / 草案/未批准规范}{来源状态：本地 PDF 已核验}
{\large\bfseries\color{Ink} 规范/Survey，无新实验；RISC-V CoVE-IO 规范草案 只支撑术语、taxonomy、接口语义或证据边界。\par}\vspace{1.8mm}
\begin{columns}[T,totalwidth=\textwidth]
  \begin{column}{0.45\textwidth}
    \fcolorbox{Rule}{Paper}{%
  \begin{minipage}[t]{0.97\linewidth}
    \vspace{1.1mm}
    \hspace{1.4mm}\begin{minipage}{0.92\linewidth}
      \PanelTitle{内容说明}
      \scriptsize \begin{itemize}
  \item 本页不展示独立系统实验，也不把二手归纳写成一手结果。
  \item 可支撑：CoVE-IO and TEE-I/O draft is the current standards-track SOTA for RISC-V confidential。
  \item 证据等级：E3 / Draft-not-ratified；来源状态：本地 PDF 已核验。
  \item 涉及具体平台、协议状态或数值时，转到原始论文、官方规范或厂商材料核验。
\end{itemize}
    \end{minipage}
    \vspace{1mm}
  \end{minipage}}
  \end{column}
  \begin{column}{0.49\textwidth}
    \fcolorbox{Rule}{Panel}{%
  \begin{minipage}[t]{0.97\linewidth}
    \vspace{1.1mm}
    \hspace{1.4mm}\begin{minipage}{0.92\linewidth}
      \PanelTitle{实验/证据边界}
      \scriptsize {\tiny
\begin{tabularx}{\linewidth}{@{}YYY@{}}
\toprule
\textbf{对象} & \textbf{可支撑} & \textbf{边界} \\
\midrule
数据/来源 & RISC-V CoVE-IO 规范草案 & 本地 PDF 已核验 \\
Baseline/对照 & 以论文或规范原文为准 & 不跨材料外推 \\
指标/对象 & 只使用当前来源可证明的信息 & E3 / 草案/未批准规范 \\
使用边界 & 可进入本方向演进图 & 不能替代更强证据 \\
\bottomrule
\end{tabularx}
}
    \end{minipage}
    \vspace{1mm}
  \end{minipage}}
  \end{column}
\end{columns}
\vspace{1mm}
{\tiny\textcolor{Muted}{引用依据：引用依据为论文原文、官方规范或公开材料；本页结论按 E3 / 草案/未批准规范 的证据等级使用，不外推到未验证平台。}}
\end{frame}


\begin{frame}[t,shrink=6]{RISC-V CoVE-IO 规范草案：文章评价}
\DeckKicker{文章评价}
\SlideMeta{证据等级}{E3 / 草案/未批准规范}{来源状态：本地 PDF 已核验}
{\large\bfseries\color{Ink} 评价：RISC-V CoVE-IO 规范草案 适合做 taxonomy/规范入口；规范/Survey，无新实验，不能替代一手系统证据。\par}\vspace{1.8mm}
\begin{columns}[T,totalwidth=\textwidth]
  \begin{column}{0.45\textwidth}
    \fcolorbox{Rule}{Paper}{%
  \begin{minipage}[t]{0.97\linewidth}
    \vspace{1.1mm}
    \hspace{1.4mm}\begin{minipage}{0.92\linewidth}
      \PanelTitle{内容说明}
      \scriptsize \begin{itemize}
  \item 设计优势：当前最重要的 RISC-V confidential I/O 规范入口。
  \item 局限性：v0.3.0 draft 仍可能变化，不证明任一硬件或 OS 实现成熟。
  \item 商业化潜力：对 RISC-V confidential accelerator/NIC 生态关键；风险在 SPDM/TDISP/IDE 与平台 ABI 收敛。
  \item 规范/Survey，无新实验；具体平台结论转到一手材料核验。
\end{itemize}
    \end{minipage}
    \vspace{1mm}
  \end{minipage}}
  \end{column}
  \begin{column}{0.49\textwidth}
    \fcolorbox{Rule}{Panel}{%
  \begin{minipage}[t]{0.97\linewidth}
    \vspace{1.1mm}
    \hspace{1.4mm}\begin{minipage}{0.92\linewidth}
      \PanelTitle{文章评价}
      \scriptsize {\tiny
\begin{tabularx}{\linewidth}{@{}YY@{}}
\toprule
\textbf{维度} & \textbf{判断} \\
\midrule
设计优势 & 当前最重要的 RISC-V confidential I/O 规范入口。 \\
主要局限 & v0.3.0 draft 仍可能变化，不证明任一硬件或 OS 实现成熟。 \\
商业落地 & 对 RISC-V confidential accelerator/NIC 生态关键；风险在 SPDM/TDISP/IDE 与平台 ABI 收敛。 \\
讲稿定位 & 代表性改进; 证据强度 E3 \\
\bottomrule
\end{tabularx}
}
    \end{minipage}
    \vspace{1mm}
  \end{minipage}}
  \end{column}
\end{columns}
\vspace{1mm}
{\tiny\textcolor{Muted}{引用依据：引用依据为论文原文、官方规范或公开材料；本页结论按 E3 / 草案/未批准规范 的证据等级使用，不外推到未验证平台。}}
\end{frame}


\begin{frame}[t,shrink=6]{RISC-V IOMMU 支撑 CoVE-IO：内容摘要}
\DeckKicker{内容摘要}
\SlideMeta{证据等级}{E0 / Spec-standard SOTA}{来源状态：本地 PDF 已核验}
{\large\bfseries\color{Ink} IOMMU 支撑设备 DMA translation/protection，是 CoVE-IO 设备侧地址访问控制的标准底座\par}\vspace{1.8mm}
\begin{columns}[T,totalwidth=\textwidth]
  \begin{column}{0.45\textwidth}
    \fcolorbox{Rule}{Paper}{%
  \begin{minipage}[t]{0.97\linewidth}
    \vspace{1.1mm}
    \hspace{1.4mm}\begin{minipage}{0.92\linewidth}
      \PanelTitle{内容说明}
      \scriptsize \begin{itemize}
  \item IOMMU 支撑设备 DMA translation/protection，是 CoVE-IO 设备侧地址访问控制的标准底座
  \item RISC-V IOMMU specification provides the current DMA isolation substrate for CoVE-IO.
\end{itemize}
    \end{minipage}
    \vspace{1mm}
  \end{minipage}}
  \end{column}
  \begin{column}{0.49\textwidth}
    \fcolorbox{Rule}{Panel}{%
  \begin{minipage}[t]{0.97\linewidth}
    \vspace{1.1mm}
    \hspace{1.4mm}\begin{minipage}{0.92\linewidth}
      \PanelTitle{证据定位}
      \scriptsize {\tiny
\begin{tabularx}{\linewidth}{@{}YY@{}}
\toprule
\textbf{字段} & \textbf{内容} \\
\midrule
论文定位 & 当前 SOTA/补充边界 \\
材料类型 & 规范/标准材料 \\
证据等级 & E0 / Spec-standard SOTA \\
来源状态 & 本地 PDF 已核验 \\
\bottomrule
\end{tabularx}
}
    \end{minipage}
    \vspace{1mm}
  \end{minipage}}
  \end{column}
\end{columns}
\vspace{1mm}
{\tiny\textcolor{Muted}{引用依据：引用依据为论文原文、官方规范或公开材料；本页结论按 E0 / Spec-standard SOTA 的证据等级使用，不外推到未验证平台。}}
\end{frame}


\begin{frame}[t,shrink=6]{RISC-V IOMMU 支撑 CoVE-IO：研究背景}
\DeckKicker{研究背景}
\SlideMeta{证据等级}{E0 / Spec-standard SOTA}{来源状态：本地 PDF 已核验}
{\large\bfseries\color{Ink} CoVE-IO 不能只靠 PMP 或 CPU page table\par}\vspace{1.8mm}
\begin{columns}[T,totalwidth=\textwidth]
  \begin{column}{0.45\textwidth}
    \fcolorbox{Rule}{Paper}{%
  \begin{minipage}[t]{0.97\linewidth}
    \vspace{1.1mm}
    \hspace{1.4mm}\begin{minipage}{0.92\linewidth}
      \PanelTitle{内容说明}
      \scriptsize \begin{itemize}
  \item CoVE-IO 不能只靠 PMP 或 CPU page table
  \item 设备侧 DMA 需要独立的 translation/protection 语义
\end{itemize}
    \end{minipage}
    \vspace{1mm}
  \end{minipage}}
  \end{column}
  \begin{column}{0.49\textwidth}
    \fcolorbox{Rule}{Panel}{%
  \begin{minipage}[t]{0.97\linewidth}
    \vspace{1.1mm}
    \hspace{1.4mm}\begin{minipage}{0.92\linewidth}
      \PanelTitle{问题边界}
      \scriptsize {\tiny
\begin{tabularx}{\linewidth}{@{}YY@{}}
\toprule
\textbf{维度} & \textbf{说明} \\
\midrule
保护对象 & RISC-V CoVE-IO / TEE-I/O \\
传统瓶颈 & 把设备身份、DMA/MMIO、interrupt、link security 与 I/O translation 放到同一边界。 \\
本文入口 & RISC-V IOMMU 支撑 CoVE-IO \\
证据限制 & E0 / Spec-standard SOTA \\
\bottomrule
\end{tabularx}
}
    \end{minipage}
    \vspace{1mm}
  \end{minipage}}
  \end{column}
\end{columns}
\vspace{1mm}
{\tiny\textcolor{Muted}{引用依据：引用依据为论文原文、官方规范或公开材料；本页结论按 E0 / Spec-standard SOTA 的证据等级使用，不外推到未验证平台。}}
\end{frame}


\begin{frame}[t,shrink=6]{RISC-V IOMMU 支撑 CoVE-IO：解决方案}
\DeckKicker{解决方案}
\SlideMeta{证据等级}{E0 / Spec-standard SOTA}{来源状态：本地 PDF 已核验}
{\large\bfseries\color{Ink} IOMMU 控制设备可访问地址空间\par}\vspace{1.8mm}
\begin{columns}[T,totalwidth=\textwidth]
  \begin{column}{0.45\textwidth}
    \fcolorbox{Rule}{Paper}{%
  \begin{minipage}[t]{0.97\linewidth}
    \vspace{1.1mm}
    \hspace{1.4mm}\begin{minipage}{0.92\linewidth}
      \PanelTitle{内容说明}
      \scriptsize \begin{itemize}
  \item IOMMU 控制设备可访问地址空间
  \item `riscv\_aia\_2023` 作为 trusted MSI/interrupt delivery 辅助材料
\end{itemize}
    \end{minipage}
    \vspace{1mm}
  \end{minipage}}
  \end{column}
  \begin{column}{0.49\textwidth}
    \fcolorbox{Rule}{Panel}{%
  \begin{minipage}[t]{0.97\linewidth}
    \vspace{1.1mm}
    \hspace{1.4mm}\begin{minipage}{0.92\linewidth}
      \PanelTitle{核心设计拆解}
      \scriptsize \begin{itemize}
  \item 01. Device DMA translation/protection
  \item 02. I/O page-table permission semantics
  \item 03. Integration point with CoVE-IO trusted device lifecycle
\end{itemize}
    \end{minipage}
    \vspace{1mm}
  \end{minipage}}
  \end{column}
\end{columns}
\vspace{1mm}
{\tiny\textcolor{Muted}{引用依据：引用依据为论文原文、官方规范或公开材料；本页结论按 E0 / Spec-standard SOTA 的证据等级使用，不外推到未验证平台。}}
\end{frame}


\begin{frame}[t,shrink=6]{RISC-V IOMMU 支撑 CoVE-IO：实验结果}
\DeckKicker{实验结果}
\SlideMeta{证据等级}{E0 / Spec-standard SOTA}{来源状态：本地 PDF 已核验}
{\large\bfseries\color{Ink} 规范/Survey，无新实验；RISC-V IOMMU 支撑 CoVE-IO 只支撑术语、taxonomy、接口语义或证据边界。\par}\vspace{1.8mm}
\begin{columns}[T,totalwidth=\textwidth]
  \begin{column}{0.45\textwidth}
    \fcolorbox{Rule}{Paper}{%
  \begin{minipage}[t]{0.97\linewidth}
    \vspace{1.1mm}
    \hspace{1.4mm}\begin{minipage}{0.92\linewidth}
      \PanelTitle{内容说明}
      \scriptsize \begin{itemize}
  \item 本页不展示独立系统实验，也不把二手归纳写成一手结果。
  \item 可支撑：RISC-V IOMMU specification provides the current DMA isolation substrate for CoVE-IO.
  \item 证据等级：E0 / Spec-standard SOTA；来源状态：本地 PDF 已核验。
  \item 涉及具体平台、协议状态或数值时，转到原始论文、官方规范或厂商材料核验。
\end{itemize}
    \end{minipage}
    \vspace{1mm}
  \end{minipage}}
  \end{column}
  \begin{column}{0.49\textwidth}
    \fcolorbox{Rule}{Panel}{%
  \begin{minipage}[t]{0.97\linewidth}
    \vspace{1.1mm}
    \hspace{1.4mm}\begin{minipage}{0.92\linewidth}
      \PanelTitle{实验/证据边界}
      \scriptsize {\tiny
\begin{tabularx}{\linewidth}{@{}YYY@{}}
\toprule
\textbf{对象} & \textbf{可支撑} & \textbf{边界} \\
\midrule
数据/来源 & RISC-V IOMMU 支撑 CoVE-IO & 本地 PDF 已核验 \\
Baseline/对照 & 以论文或规范原文为准 & 不跨材料外推 \\
指标/对象 & 只使用当前来源可证明的信息 & E0 / Spec-standard SOTA \\
使用边界 & 可进入本方向演进图 & 不能替代更强证据 \\
\bottomrule
\end{tabularx}
}
    \end{minipage}
    \vspace{1mm}
  \end{minipage}}
  \end{column}
\end{columns}
\vspace{1mm}
{\tiny\textcolor{Muted}{引用依据：引用依据为论文原文、官方规范或公开材料；本页结论按 E0 / Spec-standard SOTA 的证据等级使用，不外推到未验证平台。}}
\end{frame}


\begin{frame}[t,shrink=6]{RISC-V IOMMU 支撑 CoVE-IO：文章评价}
\DeckKicker{文章评价}
\SlideMeta{证据等级}{E0 / Spec-standard SOTA}{来源状态：本地 PDF 已核验}
{\large\bfseries\color{Ink} 评价：RISC-V IOMMU 支撑 CoVE-IO 适合做 taxonomy/规范入口；规范/Survey，无新实验，不能替代一手系统证据。\par}\vspace{1.8mm}
\begin{columns}[T,totalwidth=\textwidth]
  \begin{column}{0.45\textwidth}
    \fcolorbox{Rule}{Paper}{%
  \begin{minipage}[t]{0.97\linewidth}
    \vspace{1.1mm}
    \hspace{1.4mm}\begin{minipage}{0.92\linewidth}
      \PanelTitle{内容说明}
      \scriptsize \begin{itemize}
  \item 设计优势：CoVE-IO 中 DMA isolation 和 I/O translation 的标准支撑点。
  \item 局限性：IOMMU 不是单独的 confidential I/O 方案，也不覆盖 device identity、SPDM/TDISP 或 link security。
  \item 商业化潜力：适合服务器设备直通与多租户 DMA 管控；风险在配置完整性、固件测量和 verifier 可见性。
  \item 规范/Survey，无新实验；具体平台结论转到一手材料核验。
\end{itemize}
    \end{minipage}
    \vspace{1mm}
  \end{minipage}}
  \end{column}
  \begin{column}{0.49\textwidth}
    \fcolorbox{Rule}{Panel}{%
  \begin{minipage}[t]{0.97\linewidth}
    \vspace{1.1mm}
    \hspace{1.4mm}\begin{minipage}{0.92\linewidth}
      \PanelTitle{文章评价}
      \scriptsize {\tiny
\begin{tabularx}{\linewidth}{@{}YY@{}}
\toprule
\textbf{维度} & \textbf{判断} \\
\midrule
设计优势 & CoVE-IO 中 DMA isolation 和 I/O translation 的标准支撑点。 \\
主要局限 & IOMMU 不是单独的 confidential I/O 方案，也不覆盖 device identity、SPDM/TDISP 或 link security。 \\
商业落地 & 适合服务器设备直通与多租户 DMA 管控；风险在配置完整性、固件测量和 verifier 可见性。 \\
讲稿定位 & 当前 SOTA/补充边界; 证据强度 E0 \\
\bottomrule
\end{tabularx}
}
    \end{minipage}
    \vspace{1mm}
  \end{minipage}}
  \end{column}
\end{columns}
\vspace{1mm}
{\tiny\textcolor{Muted}{引用依据：引用依据为论文原文、官方规范或公开材料；本页结论按 E0 / Spec-standard SOTA 的证据等级使用，不外推到未验证平台。}}
\end{frame}


\begin{frame}[t,shrink=6]{RISC-V CoVE-IO / TEE-I/O：方向总结}
\DeckKicker{方向总结}
\SlideMeta{证据等级}{方向综述}{来源状态：公开材料综合}
{\large\bfseries\color{Ink} RISC-V CoVE-IO / TEE-I/O 的结论是：三篇主讲材料共同给出技术演进，但仍要用证据等级约束每个结论。\par}\vspace{1.8mm}
\begin{columns}[T,totalwidth=\textwidth]
  \begin{column}{0.45\textwidth}
    \fcolorbox{Rule}{Paper}{%
  \begin{minipage}[t]{0.97\linewidth}
    \vspace{1.1mm}
    \hspace{1.4mm}\begin{minipage}{0.92\linewidth}
      \PanelTitle{内容说明}
      \scriptsize \begin{itemize}
  \item sIOPMP: scalable I/O protection for TEE：开创/基础入口，E1 / 同行评审改进证据。
  \item RISC-V CoVE-IO 规范草案：代表性改进，E3 / 草案/未批准规范。
  \item RISC-V IOMMU 支撑 CoVE-IO：当前 SOTA/补充边界，E0 / Spec-standard SOTA。
  \item 本方向结论：先用三篇主讲材料建立演进关系，再用证据等级控制机制、实验和商业化结论。
\end{itemize}
    \end{minipage}
    \vspace{1mm}
  \end{minipage}}
  \end{column}
  \begin{column}{0.49\textwidth}
    \fcolorbox{Rule}{Panel}{%
  \begin{minipage}[t]{0.97\linewidth}
    \vspace{1.1mm}
    \hspace{1.4mm}\begin{minipage}{0.92\linewidth}
      \PanelTitle{本方向方法对比}
      \scriptsize {\tiny
\begin{tabularx}{\linewidth}{@{}YYY@{}}
\toprule
\textbf{论文} & \textbf{定位} & \textbf{选择理由} \\
\midrule
sIOPMP: scalable I/O protection for TEE & 开创/基础入口 & sIOPMP is a peer-reviewed mechanism starting point for RISC-V trusted I/O。 \\
RISC-V CoVE-IO 规范草案 & 代表性改进 & CoVE-IO and TEE-I/O draft is the current standards-track SOTA for RISC-V。 \\
RISC-V IOMMU 支撑 CoVE-IO & 当前 SOTA/补充边界 & RISC-V IOMMU specification provides the current DMA isolation substrate f。 \\
\bottomrule
\end{tabularx}
}
    \end{minipage}
    \vspace{1mm}
  \end{minipage}}
  \end{column}
\end{columns}
\vspace{1mm}
{\tiny\textcolor{Muted}{引用依据：总结页综合三篇主讲材料的选择理由、评价字段和证据等级。}}
\end{frame}
